% Created 2026-10-02 Fri 15:47
% Intended LaTeX compiler: xelatex
\documentclass[12pt]{article}
\usepackage{graphicx}
\usepackage{longtable}
\usepackage{wrapfig}
\usepackage{rotating}
\usepackage[normalem]{ulem}
\usepackage{amsmath}
\usepackage{amssymb}
\usepackage{capt-of}
\usepackage{hyperref}
\usepackage[a4paper,margin=1in]{geometry}
\usepackage{xurl}
\usepackage{graphicx}
\usepackage{parskip}
\usepackage{titling}
\makeatletter
\renewcommand{\maketitle}{%
\begin{flushleft}
\includegraphics[width=0.5\textwidth]{logo.png}
\end{flushleft}
}
\makeatother
\author{fcalado}
\date{\today}
\title{}
\hypersetup{
 pdfauthor={fcalado},
 pdftitle={},
 pdfkeywords={},
 pdfsubject={},
 pdfcreator={Emacs 29.3 (Org mode 9.6.15)}, 
 pdflang={English}}
\begin{document}

Peer Review Committee \\[0pt]
Pratt Institute, School of Information \\[0pt]
144 West 14th Street, 6th Floor \\[0pt]
New York, NY 10011-7301 \\[0pt]

Dear Members of the SI Peer Review Committee,

Kindly accept this letter and submitted materials as my application
for third-year reappointment as Assistant Professor in the School of
Information at Pratt Institute. Enclosed you will find my curriculum
vitae, as well as teaching, research, and service statements, and
copies of my course evaluations and most recent research articles.

Since joining the School of Information, I have developed a research
agenda that uses digital methods to study language, gender, sex, and
sexuality. In my teaching, I have developed courses that bring
together humanistic approaches with programming, machine learning, and
digital humanities methods. I have also contributed to Pratt through
participation to faculty committees, the Pratt Faculty Union, and to
student groups.

I will briefly highlight several accomplishments from my first two
years at Pratt:

\textbf{\textbf{Teaching}}

Since Fall 2024, I have taught five courses across the School of
Information: INFO 601: Foundations of Information; INFO 664:
Programming Cultural Heritage; INFO 657: Digital Humanities; INFO 697:
Special Topics—How to Build a Bot; and INFO 656: Machine Learning.
Across these courses, I have worked to develop a teaching practice
that brings together computational and humanistic skills, with
particular attention to close reading, writing, and the contingent
nature of knowledge.

In Summer 2025, I taught a new course, INFO 697: How to Build a Bot,
which combines hands-on programming with critical examination of
coding, automation, and the social implications of bots. I plan to
continue developing this area of the curriculum through a Python
course focused specifically on working with text-based data, building
on a workshop I currently offer through the Data Storytelling Lab.

In addition to classroom teaching, I currently advise approximately 35
student, discussing coursework, projects, portfolios, and career
goals. I also advise the \emph{Stacks} Student Journal and Latines in Tech
student groups. These activities have allowed me to extend my teaching
and mentorship beyond individual courses and to contribute to
students' broader intellectual and professional development.

\textbf{\textbf{Research}}

My research since joining Pratt has established a sustained
interdisciplinary agenda connecting queer and trans studies, digital
humanities, and AI ethics. As a language scholar trained primarily in
English literature, I use computational tools to examine how language
related to gender, sex, and sexuality is represented, processed, and
transformed through information technologies. A central concern of
this work is understanding how computational mechanisms shape what can
be seen, represented, and said in language data.

From Fall 2024 through Fall 2026, I have written five sole-author
papers, two of which have been published in peer-reviewed venues and
three of which are currently under review. Additionally, my book
project, \emph{Queer Code: Re-working Language Technologies}, is under
contract with Bloomsbury Academic Press, with the manuscript due in
March 2027.

Across these projects, I have sought to communicate across fields that
do not always share methods or assumptions. This interdisciplinary
orientation is central to my role as a faculty member in the School of
Information, where I approach information technologies as systems that
shape how knowledge, categories, and meanings become visible.

\textbf{\textbf{Service}}

In addition to my research and teaching, I have contributed to the
School of Information through student mentoring and advising and
through service to Pratt and broader scholarly communities. I have
also participated in the Pratt Faculty Union and external scholarly
communities, as detailed in my curriculum vitae and accompanying
materials.

Taken together, my work has allowed me to build exciting momentum
across teaching, research, and service areas at Pratt SI. I am
grateful for the opportunity to contribute to the school, and look
forward to furthering the work I've begun as a continuing member of
the Pratt full-time faculty.

Thank you for your consideration.

Sincerely,

Filipa Calado \\[0pt]
Assistant Professor \\[0pt]
School of Information \\[0pt]
Pratt Institute
\end{document}
